% =====================================================================
% Green IT 2026 (MICS2-63): Short Paper, HPC Job Track
%
% This is the document you write if you assess an HPC JOB.
% If you assess a web page, write src/short-paper-mics.tex instead and
% delete this file.
% It compiles as shipped.
%
% Length: minimum 4 double-column pages, EXCLUDING guidance text,
% references and appendices. Printed on page 1 under the author block, so
% that a student who opens the PDF sees the rule.
%
% Build:  latexmk        (run it from the folder above this one)
% =====================================================================
\documentclass[conference,a4paper]{IEEEtran}

\usepackage[T1]{fontenc}
\usepackage{graphicx}
\usepackage{booktabs}
\usepackage{url}
\usepackage{amsmath}

% ---------------------------------------------------------------------
% Guidance notes. Each one tells you what its section owes.
% Delete a note as you write its section, and set \guidancetrue below to
% \guidancefalse before you submit: the PDF you hand in carries no guidance.
% The notes add pages, so that switch is also how you measure your real length.
\newif\ifguidance
\guidancefalse
\newcommand{\guidance}[1]{\ifguidance{\itshape\small #1\par\medskip}\fi}

\begin{document}

\title{Life Cycle Assessment of Exact 31-mer Counting on an HPC Node: Carbon Payback of an AI-Assisted Optimisation}

\author{\IEEEauthorblockN{Isaac Palma-Medina}
\IEEEauthorblockA{Master in High Performance Computing\\
University of Luxembourg\\
isaac.palma.001@student.uni.lu\\[2pt]
% The page minimum, printed here so it is full width and survives \guidancefalse.
\small\itshape Minimum of 4 double-column pages, excluding guidance text, references and appendices.\normalfont\upshape\\[2pt]
% Build stamp: the date this file was last compiled.
\scriptsize Compiled \today.}}

\maketitle


\begin{abstract}
\guidance{Maximum 200 words, one paragraph. State the job assessed, the
research question, what you did, the headline numbers (baseline
gCO\textsubscript{2}e per run, post-optimisation gCO\textsubscript{2}e per run,
payback in runs), \textbf{all three measured the same way and each named with
the partition and the allocation shape they were measured on}, and what
generalises beyond your job. Every sentence states a result.}
\end{abstract}

\guidance{\textbf{This is the HPC job track.} The archive also holds
\texttt{src/short-paper-mics.tex}, which is the same paper written for a web
page. Write one of the two. Delete the other, and the build follows.}

\guidance{Every italic note in this document is guidance.
Delete each one as you write its section, then set
\texttt{\textbackslash guidancefalse} in the preamble: the PDF you submit
carries no guidance. The notes add pages, so that switch is also how you check
whether you have reached the required length.}

\guidance{Keep the first half of the title above, which names the job you
assessed. Adapt the second half to your own finding: once you hold your payback
figure, rewrite it around that. A title that states a finding carries further
than one that states a topic.}

\guidance{Your paper is reviewed double-blind. Two classmates receive it with
the author block blanked, with Appendix B removed, and with the repository
replaced by an archive holding only the script, its README and the LICENSE
with the name neutralised. Your reviewers may be working on a web page rather
than on a job, so every convention your paper runs on is stated in the paper.}

% ---------------------------------------------------------------------

\section{Introduction}\label{sec:intro}
\guidance{Context, the assessed job, and the research question. The
research question is about the carbon payback of AI-assisted
optimisation, anchored on your own job. State it explicitly, in one
sentence, ending with a question mark.}

% ---------------------------------------------------------------------
\section{Related Work}
\guidance{Between 3 and 5 references, \emph{all of them found by you}.
Finding the work that surrounds your question is part of the exercise,
and it is where an AI assistant needs checking most.

Position your question with respect to them. Do not produce a list of
summaries: each reference must earn its place by supporting, framing or
contradicting something you do.

Every reference must be verified by you, by hand. Any reference an AI
proposed that turns out not to exist must be recorded in Appendix B. A
reference that does not exist costs you this section.

\emph{What is expected in v1, and what in v2.} The method for searching is
taught in week 3, six weeks before v1 is due, so you have time to read a
field. \textbf{In v1 this section is marked on three things: your references
are there, every one is verified, and the section situates what you did
against them}. Your reviewers will tell you what is missing, and
\textbf{v2 is where that section earns its keep}.}

% ---------------------------------------------------------------------
\section{Life Cycle Assessment}
\guidance{The five subsections below follow Session 3 with respect to
ISO 14040~\cite{iso14040} and ISO 14044~\cite{iso14044}. In each subsection,
introduce its intention concisely, then apply it to your job.}

\subsection{Goal and scope definition}\label{sec:goalscope}
The goal and scope definition fixes what is studied, the functional unit every
later figure refers to, the boundary of the system and the method that turns
inventory flows into an impact~\cite{iso14040,iso14044}.

\emph{Goal.} This study quantifies the carbon footprint of one execution of a
baseline 31-mer counting job. That footprint is the reference against which
the AI-assisted optimisation of Section~\ref{sec:optimisation} and its carbon
payback in Section~\ref{sec:payback} are measured.

\emph{The job.} The code is \texttt{kc-py1.py} from the public repository
\texttt{lh3/kmer-cnt}~\cite{kmercnt}, frozen at commit
\texttt{e2574719cfb7\allowbreak 84915d80eb58\allowbreak 28e78dfae4cf\allowbreak dd7b}.
It is a single-threaded Python~3 program. It reads FASTA from standard input and
counts every canonical 31-mer, the lexicographically smaller of a 31-mer and its
reverse complement. It skips any 31-mer that contains \texttt{N}. It prints a
histogram with 255 lines, one per count from 1 to 255, and pools every count
above 255 into the last line. The input is
\texttt{M\_abscessus\_HiSeq\_10M.fa.gz} from the repository's v0.1
release~\cite{kmercnt}: 187\,693\,955~bytes, SHA-256
\texttt{dd2ef49aebab\allowbreak fed1574b55fd\allowbreak 002025f29b2f\allowbreak da88c70b2095\allowbreak f5e1b3a822ce\allowbreak e736}.
It is decompressed on the fly by \texttt{gzip}. The program upper-cases only the
last record of its input. No sequence line of this input holds a lower-case
base (checked on 28~September~2026), so this does not affect the result. The job
is submitted with the Slurm script \texttt{run-baseline.sh} of the project
repository at commit
\texttt{253d93ec726b\allowbreak 82670c8f4b84\allowbreak da20e9528c9a\allowbreak 59ef}.
Every run writes the same 1\,972-byte histogram, SHA-256
\texttt{06a32de77200\allowbreak 121a397050e0\allowbreak 15a257953dc9\allowbreak 56b7854027a1\allowbreak ed73b15afc12\allowbreak bf42}.

\emph{The machine.} The job runs on the Aion cluster of the University of
Luxembourg HPC platform, partition \texttt{batch}, QoS \texttt{normal}, on the
regular compute node \texttt{aion-0127}. An Aion node holds two AMD EPYC~7H12
processors of 64~cores each, 256~GB of DDR4 memory, a 480~GB local SSD and one
HDR100 InfiniBand port~\cite{ulhpcaion}. The loaded modules are
\texttt{env/release/default} and \texttt{lang/Python/3.11.5-GCCcore-13.2.0}
(Python~3.11.5), as recorded in every run's log.

\emph{The allocation.} Each run holds one node exclusively: one task,
128~cores, all the memory Slurm can reserve on the node
(\texttt{-{}-mem=0}, recorded as 224\,000~MiB) and a walltime limit of
30~minutes. The program uses one of those cores. The node is reserved
exclusively so that no other job shares it. As a consequence, the whole node,
not one core, is attributed to the job.

\emph{The user.} The job is submitted with \texttt{sbatch} from an Aion login
node by the author, through a standard student account.

\emph{How many runs.} The functional unit is one run. Five runs are measured
(Section~\ref{sec:lci}). The job is assumed to be executed 100~times per year,
an assumption listed in Section~\ref{sec:threats}.

\emph{Functional unit.} The GWP100 carbon footprint of one execution of
\texttt{kc-py1.py} at commit \texttt{e2574719} of \texttt{lh3/kmer-cnt},
counting exact canonical 31-mers in \texttt{M\_abscessus\_HiSeq\_10M.fa.gz}
(187\,693\,955~bytes, SHA-256 \texttt{dd2ef49a\ldots e736}), submitted by a
student user on Aion, partition \texttt{batch}, holding one exclusive regular
node (2~$\times$~AMD EPYC~7H12, 128~cores, 224\,000~MiB reservable memory) for
a median walltime of 368~s.

\emph{Method.} One criterion: a carbon footprint in CO\textsubscript{2}e.
The characterisation method is IPCC~2021~\cite{ipcc2021} and the indicator is
GWP100.

\emph{Boundary.} Inside the boundary are the node the job holds for its
walltime, and the share of the interconnect, the storage and the site that
serve it. The last three enter through the power usage effectiveness (PUE)
declared in Section~\ref{sec:lci}. Outside the boundary are the author's laptop,
the login node, the network between laptop and cluster, the one-off download of
the input, the development of the code, and every node and job the allocation
does not reserve.

\emph{Life cycle phases.} For the compute node, the study keeps production,
including the extraction of raw materials that cradle-to-gate data contain, and
use. For the interconnect, the storage and the site, it keeps use only, through
the declared PUE. End of life is not inventoried separately: the disposal route
of Aion's hardware is not published. One of the candidate hardware datasets
of Section~\ref{sec:lci} nonetheless includes disposal. That mismatch is listed
in Section~\ref{sec:threats}.

\emph{Components.} The study is capped at four components, which
Section~\ref{sec:lci} inventories in this order:
(1)~the \emph{compute node}: its electricity during the walltime and a share of
its manufacture;
(2)~the \emph{interconnect}: the HDR100 InfiniBand fabric that carries the
job's reads and writes between the node and shared storage;
(3)~the \emph{storage}: the shared GPFS file system that holds the input, the
code and the output;
(4)~the \emph{site}: cooling and power distribution around the machine.

\subsection{Life cycle inventory}\label{sec:lci}
The inventory lists the physical flows the assessment runs on, and how each
was obtained. It holds no grams: those belong to Section~\ref{sec:lcia}. Every
entry is gathered in Table~\ref{tab:inventory} in Appendix~A, with its source
and the date it was read. The raw scheduler records, logs and commands are kept
in the project repository.

\emph{The measurement.} The frozen commit and input of
Section~\ref{sec:goalscope} were run five times as separate batch jobs on
29~September~2026, submitted consecutively with
\texttt{-{}-dependency=singleton} so that no two overlapped. Each run held its
node exclusively, and all five landed on \texttt{aion-0127}. All five completed
with exit code~0 and wrote identical output. All five are kept and the median
is reported. Walltimes were 362, 364, 368, 369 and 372~s: a median of 368~s and
a spread of 10~s from fastest to slowest. File system caches were neither reset
nor measured. The first run is not the slowest (the third is), so these runs
show no first-run effect, but they do not establish whether any run started
cold. The records were read with \texttt{sacct} on 29~September~2026. The
fields read, and why each is the quantity claimed, are:
\texttt{Elapsed}, the time from start to end of the allocation, which is the
walltime the node is held;
\texttt{AllocNodes} and \texttt{AllocCPUS}, the node and cores held;
\texttt{ReqMem}, the memory reserved;
\texttt{MaxRSS} of the Python step, the peak resident memory the program
touched, which agrees with the maximum resident set size that
\texttt{/usr/bin/time -v} reports to within 0.2\,\%;
\texttt{MaxDiskRead} and \texttt{MaxDiskWrite} of the batch and Python steps;
and \texttt{ConsumedEnergy} and \texttt{ConsumedEnergyRaw}.

\emph{Whether an energy figure comes back.} It does not come back from the
scheduler. On 29~September~2026, \texttt{scontrol show config} reported the
energy plugin \texttt{acct\_gather\_energy/rapl}. Nevertheless, both energy fields
were empty for every step of all five jobs, and for a separate interactive
allocation checked the same day. The empty fields are missing data, not zero
energy. The processor's own RAPL counters are readable by an ordinary user
through \texttt{/sys/class/powercap}, while \texttt{perf} reports its package
energy event as not supported. The job was therefore run five more times on
29~September~2026, with the same program, input and allocation. Each run read
both package counters directly before and after the Python step, and corrected
for counter wrap. \texttt{aion-0127} was busy, so these five runs and an idle
reading were moved to \texttt{aion-0144}, an identical regular node. All five
completed with exit code~0 and wrote the same histogram, and Slurm's energy
fields stayed empty. Over Python steps of 357.6 to 366.4~s, the two CPU packages
drew a median of 67\,834~J (range 67\,129 to 68\,827~J), an average of 188~W.
Read the same way for 60~s after the fifth run, the idle node drew 10\,959~J, or
183~W. Assuming that idle rate held throughout each run, the job's activity
adds a median of only 1\,863~J (range 1\,690 to 1\,953~J). So 97\,\% of the
package energy is drawn by a node the job holds but barely uses. RAPL covers
the processor packages only: memory, board, local disk, fans and network card
are outside it, and so is the time the allocation runs outside the Python step.
Because no instrument reaches the whole node, the node's electricity is also
estimated with the Green Algorithms model~\cite{lannelongue2021}:
the walltime times the power of the cores and memory held, using 280~W per
64-core processor~\cite{ulhpcaion} and 0.3725~W per GB of memory. The whole
256~GB installed is counted, because the exclusive allocation holds all of it.
Counting one active core gives 36.7~kJ (0.0102~kWh) per run. Counting all
128~allocated cores at full power gives 241~kJ (0.0670~kWh). The measured
package energy alone, 67.8~kJ, is almost twice the first figure. So the one-core
model understates the node, and both model figures are scenarios, not bounds.
The measured package energy is the figure this inventory carries forward.

\emph{Allocation rule.} Anything shared across many runs is divided down to one
run by the node-seconds the job holds over the node-seconds the hardware
delivers during its service life. For one exclusive node this is the walltime
over the service life. The same rule is applied in every subsection below.

\subsubsection{Compute node}
Operational: 67\,834~J of CPU-package electricity per run (median, RAPL), of
which 1\,863~J is above idle, over a median walltime of 368~s,
holding 1~node and 128~cores. Memory reserved: 224\,000~MiB (\texttt{ReqMem}),
of the 256~GB installed~\cite{ulhpcaion}. Memory touched: 7\,760\,436~KiB, or
7.40~GiB (median \texttt{MaxRSS}). The job therefore touches 3.4\,\% of what it
reserves, and 211~GiB stays reserved but unused.
Embodied: ecoinvent cut-off v3.11~\cite{ecoinvent311} holds no server
production activity. Searches on 29~September~2026 for \emph{server} (58
activities, 32 products), \emph{computer production}, \emph{computer, desktop}
and \emph{computer, laptop} found none. The closest dataset is
\emph{computer production, desktop, without screen} (GLO, 1998--2024,
per unit). It models a 2002 single-processor desktop of 11.3~kg and includes
packaging and disposal. Under IPCC~2021 GWP100 it carries 221~kg
CO\textsubscript{2}e per unit. A two-socket AMD EPYC server, the Dell PowerEdge R7525, has a published production footprint of
1\,768.4~kg CO\textsubscript{2}e, of which 854.7~kg comes from two 4~TB
SSDs~\cite{dellr7525}. That study uses IPCC~2013 rather than IPCC~2021, and its
server differs from an Aion blade. Neither figure describes an Aion node. Which
one enters Section~\ref{sec:lcia}, and why, is decided there. Aion entered
production in October--November 2021 and still runs in September
2026~\cite{ulhpcaion}. No retirement date is published, so a service life of
six years is assumed (Section~\ref{sec:threats}). One run then holds 368~s of
189\,345\,600~node-seconds, a share of $1.94\times10^{-6}$ of the node.

\subsubsection{Interconnect}
The job runs on one node, so no node-to-node traffic occurs. Its only traffic
on the HDR100 fabric (link rate~100, read with \texttt{ibstat}) is file system
I/O. The step counters do not measure storage traffic directly. The Python
step's \texttt{MaxDiskRead} of 680.67~MiB equals the decompressed stream that
\texttt{gzip} writes into the pipe (713\,674\,148~bytes, or 680.61~MiB). The
batch step's \texttt{MaxDiskWrite} of 680.72~MiB is \texttt{gzip} writing that
same pipe. What the batch step reads beyond the pipe, about 198~MiB, is the
179~MiB compressed input plus about 19~MiB of other files, most likely the
loaded modules. The data read from storage per run is therefore about
187\,693\,955~bytes plus about 19~MiB, and the data written is the
1\,972-byte histogram plus the job log. This breakdown is inferred from the
counters, not measured on the fabric. Runs~2 to~5 may have read the input from
the node's page cache. The fabric's energy is not metered per job. It enters
through the declared PUE of the site subsection.

\subsubsection{Storage}
The repository, input and output sit on Aion's shared GPFS file system
(\texttt{df -T}, 29~September~2026)~\cite{ulhpcaion}. Per run, the job keeps
the 187\,693\,955-byte input and writes the 1\,972-byte output, reading and
writing the bytes listed under the interconnect. Storage energy is not metered
per job. It enters through the declared PUE of the site subsection.

\subsubsection{Site}
The University of Luxembourg HPC centre publishes no PUE for Aion. This study
therefore declares 1.36, the EU average PUE for 2023, reported in 2024 and
published in the European Commission's first technical report on data-centre
energy performance~\cite{eudc2025}, read on 29~September~2026. It is a regional proxy,
not a measurement of this site. It multiplies the operational energy of the
compute node, and covers cooling, power distribution, interconnect and storage.
How far a different value would move the result is discussed in
Section~\ref{sec:threats}.

\subsection{Life cycle impact assessment}\label{sec:lcia}
\guidance{Single-criteria carbon footprint: climate change only. Say so,
and say what that leaves out.

Name the characterisation method, IPCC 2021, and the indicator you read out of
it, GWP100. Hold that one indicator across every factor in the paper.

Then name every factor, where it comes from, the date you read it \emph{and its
basis}: the grid intensity of the country the machine sits in, the PUE you
declare with the argument that supports it, and then ecoinvent or a vendor
product carbon footprint sheet for anything made. A local grid factor
substitutes into the \emph{operational} term only, since hardware is made
worldwide.

Then the result: your baseline in gCO\textsubscript{2}e per run, with the units
carried at every step so a reader can follow the chain from joules to grams.
Your inputs carry two significant digits, so report your answer to two.}

\subsection{Interpretation}\label{sec:interpretation}
\guidance{What that number is a number \emph{of}, in one sentence.

Then make it mean something. Put it next to the everyday comparators of
Session 2, a kilometre by car, an hour of a 5\,W LED, a day of a Luxembourg
resident, and say what one run of your job amounts to in each. Pick the
comparison that a reader outside this course would understand first.

Then cross-check your own arithmetic. Estimate the energy from the top down,
the nodes you held times their rated power times the walltime, and put it next
to the figure Section~\ref{sec:lci} settled on for the same run. Say by how much
the two differ, and which one your result uses. Say also which kind of check it
is: a top-down estimate next to a counter reading tests the reading, and next to
a substitute it tests the substitute.}

\subsection{Reporting}\label{sec:reporting}
\guidance{Two short paragraphs, written for a reader who did not follow
the calculation.

First, the result: the number, the functional unit it belongs to, and the
conventions a reader needs in order to reuse it.

Then \emph{the hotspots}. The rest of the paper hangs on them. Which component,
or which resource, carries most of your number? Name one, two at most, with the
share each takes. \textbf{What you name here is what
Section~\ref{sec:optimisation} attacks.}}

% ---------------------------------------------------------------------
\section{Optimisation}\label{sec:optimisation}

\subsection{Optimisation procedure}\label{sec:optproc}
\guidance{Start from the hotspot you named in
Section~\ref{sec:reporting}. Say how you propose to remove it \emph{in your
job}: what you change, and why that attacks that hotspot rather than something
else. Reserved memory nobody uses, a serial phase that holds every node idle, a
file rewritten at every iteration and an allocation larger than the code can
fill are the usual four.

Then the mechanics: the commit you start from, the input you hold fixed, and
the partition and allocation shape you run on. \textbf{Before you change
anything, measure that starting point with the protocol of
Section~\ref{sec:lci} and convert it with the factors of
Section~\ref{sec:lcia}}: that pair, the quantities and the grams that follow
from them, is your before.
Then say which optimisations the AI coding agent applied, what you asked for
and what you accepted, because you are the author of both.}

\subsection{Implementation}
\guidance{What actually changed in the code or in the job script, concretely
enough that a reader could make the same changes. Anything the agent proposed
and you rejected belongs here too, with why. \textbf{Say how you checked that
the job still produces the same result}: an optimisation that changes the output
is a different job, and its numbers are no longer comparable.}

\subsection{Footprint after optimisation}\label{sec:afteropt}
\guidance{Go through Sections~\ref{sec:lci} to~\ref{sec:reporting} again, this
time on the optimised job, and measure it the same way you measured the
original. Run it on the same partition and the same node type, with the same
fixed input, and keep exactly the same functional unit as in
Section~\ref{sec:goalscope}. Those rules are what make the before and the after
comparable: a different node type changes the energy per core-second, and your
saving with it.

Then report only what the optimisation changed: the quantities that moved, the
new impact figures that follow from them, what the new result means, and the
hotspots if they are no longer in the same order. If something did not change,
say so in one line. Finish by saying what you changed to produce the
difference.}

\subsection{What generalises}
\guidance{This is the part of your work that outlives your job. It
is two things, in this order.

\textbf{1. The method.} You have just fixed one hotspot in one job. Now
write that fix as a procedure for \emph{any} job. Ordinary prose is fine,
but give it a shape: how a reader recognises the same problem in their own
job, what they should change, and how they check that it worked. Somebody
holding different code must be able to follow it without ever seeing
yours. If a step only makes sense for your job, it belongs in
Section~\ref{sec:optproc}.

\textbf{2. The tool.} A script, built with AI assistance, that carries out
that method. It must run on a job it has never seen, given nothing but its
identifier. Give the repository
URL and the commit hash, and show that it reproduces the figure you
reported. \textbf{Your script sits with a \texttt{README} of its own}: what it
does, how to run it, and what it prints. The repository you cloned this
template into already carries a \texttt{README} about building the paper, and
that one answers a different question. Your reviewer receives it as an archive
and runs it on a machine that has never seen your project, and possibly with no
cluster at all, so say what it needs and give it something to read when the
job record is out of reach.

The method to write up is the one you actually worked out on your own
job: that is the one you can defend, and the one that will still be
true in January.}

% ---------------------------------------------------------------------
\section{Carbon Payback}\label{sec:payback}
\guidance{Divide the full carbon cost of producing your optimisation by the
saving per run you measured in Section~\ref{sec:afteropt}. Give the result in
runs, then in a period, using how often the job is really submitted.

\textbf{The numerator is the cost of producing the optimisation}, and
Appendix~B already computes most of it: the estimated gCO\textsubscript{2}e of
the AI assistance that produced Section~\ref{sec:optimisation}, plus the runs
you spent measuring and testing, which on a cluster are themselves an
allocation. Count those runs with the same factors as your baseline, and say
which of the two parts you kept.

\textbf{Check that the two numbers you divide come from the same machine.} A
before measured on one partition and an after measured on another are not a
saving, they are a comparison of two node types. If the queue moved you, rerun
the before.}

\subsection{What this means}
\guidance{Answer the question you asked in Section~\ref{sec:intro}, in one
sentence, with the number. Then what a reader should take away, and what you
would do next. Keep to results you have already reported, and to claims your
figures carry.}

% ---------------------------------------------------------------------
\section{Threats to Validity}\label{sec:threats}
\guidance{The threats you could not resolve, and how far each one could move
your number. Data gaps, estimation choices, and anything the functional unit
excludes.

Six that this course hands you, and that belong here:

\emph{Every substitute you could not source.} Section~\ref{sec:lci} carries the
substitutes you could source. The ones you could not are assumptions, and each
one is named here with what it stands in for and how far it could move your
result.

\emph{What your energy figure covers.} A counter reading is taken from
the nodes your job held, and it describes your job alone when the allocation is
exclusive, which is why your protocol asks for one. On a shared node it
also carries whatever ran beside you. Whichever figure Section~\ref{sec:lci}
settled on, a reading or a substitute, the switches that carried the job's
traffic, the storage that served its files and the cooling that kept it running
sit outside it, which is why the site enters your study through a declared PUE
rather than through a meter. Say what that leaves uncertain, and if you
substituted, how far the substitute could move your result.

\emph{The PUE you declared}, its argument and its source, and what a different
value would do to your result.

\emph{Figures that are not at the same level of completeness.} One of your factors may be cradle to
grave while another covers use only. Adding them treats them as if they were the
same kind of thing. That is often the only data you have. Name every figure
whose boundary differs from the others, and say in which direction it pushes
your result.

\emph{Burden shifting}: what a carbon-only criterion hides for your job. An
improvement that lowers your number can raise a total you did not measure, and a
job made faster is often a job run more often.

\emph{The limit of one criterion}: these results cover climate change
only, on one job, and support no comparison with anybody else's job.}

% ---------------------------------------------------------------------
% Unnumbered on purpose, for two reasons. The paper has six numbered
% sections, named in Session 1 and in the marking grid, so a numbered
% seventh would shift every reference to them. And a disclosure statement
% sits outside the argument anyway, which is where journals place it.
% It exists because Appendix B is removed at distribution: without this
% paragraph, a reviewer reads a paper whose AI assistance is declared
% nowhere.
\section*{AI Usage Transparency}
\guidance{Three or four sentences, and the only place your reviewers read
about your AI use.

Name the models you worked with, say what you used each one for and which
sections they touched. State
that you verified every output you kept, and that you are the author of this
paper and of the script of Section~\ref{sec:optimisation}. Close with the
total estimated gCO\textsubscript{2}e of that assistance, the figure
Appendix~B computes.

Two things stay out of it. Anything that points at you: no file path, no
repository URL, no folder name, since this paragraph travels with the
anonymised copy your reviewers receive. And the detail: the per-model
counts, the arithmetic and the wrong answers you caught belong in
Appendix~B.}

\ifguidance
\noindent\textit{Worked example, written from the two journal entries of
Appendix~B.}

Two models assisted this work. ChatGPT (GPT-5.6 Sol) was used in chat, to
explain the inventory and impact concepts I had not understood and to review
the prose of this paper. DeepSeek V4-Flash was used agentically, to apply the
optimisations to the job, to write the script that carries the method to any
job, and to re-run the measurement protocol after each change. Every
output I kept was checked against a source I read myself, and the six wrong
answers I caught are listed in Appendix~B. I am the author of this paper and
of the script. The assistance amounts to 251\,000 output tokens, estimated at
28 gCO\textsubscript{2}e in Appendix~B.
\fi

% ---------------------------------------------------------------------
\bibliographystyle{IEEEtran}
\bibliography{refs}

% ---------------------------------------------------------------------
\guidance{\medskip
\noindent\textbf{Appendices.}
\textbf{Every figure and every table in your paper goes in
Appendix~A.} Do not put a single figure or table in the body: refer to them
from the text, and place them all there.}

\appendices
\section{Figures and Tables}
\guidance{Every figure and every table of your paper.}

\begin{table*}[!t]
\caption{Life cycle inventory of one run (Section~\ref{sec:lci}). Measured
rows are the median of five runs on \texttt{aion-0127}, except the RAPL rows
(\texttt{aion-0144}); all rows read on 29~September~2026 unless stated.}
\label{tab:inventory}
\centering\footnotesize\setlength{\tabcolsep}{4pt}
\begin{tabular}{@{}lllll@{}}
\toprule
Component & Quantity & Value and unit & Source & Kind \\
\midrule
Compute node & Walltime & 368~s (range 362--372) & \texttt{sacct} \texttt{Elapsed}, jobs 15860349--53 & measured \\
             & Nodes, cores held & 1 node, 128 cores & \texttt{sacct} \texttt{AllocNodes}, \texttt{AllocCPUS} & measured \\
             & Memory reserved & 224\,000~MiB & \texttt{sacct} \texttt{ReqMem} & measured \\
             & Memory touched & 7\,760\,436~KiB (7.40~GiB) & \texttt{sacct} \texttt{MaxRSS}, step \texttt{.0} & measured \\
             & Electricity, scheduler & none returned & \texttt{ConsumedEnergy(Raw)} empty & measured \\
             & CPU-package energy per run & 67\,834~J (range 67\,129--68\,827) & RAPL, jobs 15861652--60 & measured \\
             & CPU-package power, idle & 183~W (10\,959~J in 60~s) & RAPL, job 15861674 & measured \\
             & Package energy above idle & 1\,863~J (range 1\,690--1\,953) & run $-$ idle rate $\times$ step time & derived \\
             & Electricity, 1 core active & 36.7~kJ & Green Algorithms~\cite{lannelongue2021} & model \\
             & Electricity, 128 cores at TDP & 241~kJ & Green Algorithms~\cite{lannelongue2021} & model \\
             & Hardware & 2 $\times$ EPYC 7H12, 256~GB & \cite{ulhpcaion}, \texttt{lscpu} & looked up \\
             & Manufacture, ecoinvent proxy & 221~kg CO\textsubscript{2}e per unit & \cite{ecoinvent311}, desktop without screen & looked up \\
             & Manufacture, vendor proxy & 1\,768.4~kg CO\textsubscript{2}e per unit & \cite{dellr7525}, R7525, IPCC 2013 & looked up \\
             & Service life & 6~years & Section~\ref{sec:threats} & assumption \\
             & Share per run & $1.94\times10^{-6}$ & 368~s / 189\,345\,600~s & derived \\
Interconnect & Bytes read from storage & $\approx$187.7~MB + 19~MiB & \texttt{MaxDiskRead}, batch minus pipe & inferred \\
             & Bytes written to storage & 1\,972~B + log & \texttt{stat} of output & measured \\
Storage      & Input stored & 187\,693\,955~B & \texttt{stat}, SHA-256 checked (28 Sep) & measured \\
             & Output stored & 1\,972~B & \texttt{stat} & measured \\
Site         & PUE & 1.36 & EU average, 2023 data~\cite{eudc2025} & declared proxy \\
All          & Runs per year & 100 & Section~\ref{sec:threats} & assumption \\
\bottomrule
\end{tabular}
\end{table*}

% The AI Usage Journal starts on a page of its own. \clearpage rather than
% \newpage: in a two-column class \newpage ends the current column, while
% \clearpage ends the whole page and flushes every float still waiting, so
% a figure belonging to Appendix A is placed before the break.
\clearpage

\section{AI Usage Journal}
\guidance{For every AI model used: the model and its version, the total number
of prompts, the total input and output tokens, the estimated
gCO\textsubscript{2}e and how you got there, and a comparison with a familiar
activity. Convert with the course energy factor given in Session~7, applied to
your \emph{output} tokens, and the grid factor you declared in
Section~\ref{sec:lcia}. Cover every kind of AI use, chat assistants and agentic
tools alike.

Then, in prose: the kinds of task you asked the model for \emph{and which
sections of this paper they touched}, and a numbered list of \emph{every} piece
of wrong information it gave you, with how you spotted each one.

Naming the sections is what the University's assessment directive asks for: a
graded task must indicate, in a footnote or an appendix, which parts of it an AI
generated or influenced.

\textbf{The total you reach here is the numerator of
Section~\ref{sec:payback}}, so compute it before you write that section.

Copy the empty block below once per model and fill it in. The two entries that
follow it are worked examples. They disappear with the guidance.}

% Blank entry: hidden until filled in, then move it out of the \ifguidance.
\ifguidance
\noindent\textbf{AI model:} <model and version> \\
\textbf{Interaction:} <chat or agentic> \\
\textbf{Prompts:} <n> \\
\textbf{Input tokens:} <n> \quad \textbf{Output tokens:} <n> \\
\textbf{Energy factor:} 1.03 Wh per 1000 output tokens \\
\textbf{Grid factor:} <value> gCO\textsubscript{2}e per kWh, <source>, read on
<date> \\
\textbf{Estimated:} <value> gCO\textsubscript{2}e \\
\textbf{Comparable to:} <value> hours of a 5\,W LED, or <value> km in a small
petrol car

\textbf{Kinds of task, and which sections they touched:} <...>

\textbf{Wrong information:} <numbered list, or ``none''>


\noindent\textbf{AI model:} ChatGPT, GPT-5.6 Sol \\
\textbf{Interaction:} chat \\
\textbf{Prompts:} 74 \\
\textbf{Input tokens:} 96\,400 \quad \textbf{Output tokens:} 38\,200 \\
\textbf{Energy factor:} 1.03 Wh per 1000 output tokens \\
\textbf{Grid factor:} <value> gCO\textsubscript{2}e per kWh, <source>, read on
<date> \\
\textbf{Estimated:} <value> gCO\textsubscript{2}e \\
\textbf{Comparable to:} <value> hours of a 5\,W LED, or <value> km in a small
petrol car

\textbf{Kinds of task, and which sections they touched:} teaching me the concepts
I had not understood, such as what the inventory holds and what the impact
assessment adds to it (III.B, III.C); reviewing the prose of the whole paper;
checking the arithmetic of III.C and V.

\textbf{Wrong information:}
\begin{enumerate}
\item it proposed ``the functional unit is the code'', which is wrong: it is one
  run of that code under the conditions of Section~\ref{sec:goalscope}. I kept my
  own sentence.
\item it gave a Luxembourg grid factor with no source. I replaced it with the
  one I had read myself. I checked the basis.
\item it offered a reference with the DOI \texttt{10.1109/JXXQ.2024.193422}. I
  checked it on doi.org, it did not resolve, and the publication does not
  exist. I dropped the reference.
\end{enumerate}

\medskip
\noindent\textbf{AI model:} DeepSeek V4-Flash, pinned as
\texttt{deepseek/deepseek-v4-flash-0731} \\
\textbf{Interaction:} agentic, through OpenCode with the course OpenRouter key \\
\textbf{Prompts:} 212 requests over 9 sessions \\
\textbf{Input tokens:} 5\,140\,000 \quad \textbf{Output tokens:} 213\,000 \\
\textbf{Energy factor:} 1.03 Wh per 1000 output tokens \\
\textbf{Grid factor:} <value> gCO\textsubscript{2}e per kWh, <source>, read on
<date> \\
\textbf{Estimated:} <value> gCO\textsubscript{2}e \\
\textbf{Comparable to:} <value> hours of a 5\,W LED, or <value> km in a small
petrol car

\textbf{Kinds of task, and which sections they touched:} applying the
optimisations of Section~\ref{sec:optimisation} to the job (IV.B); writing the
script that carries the method to any job (IV.D); re-running the measurement
protocol after each change (IV.C).

\textbf{Wrong information:}
\begin{enumerate}
\item it reported that the job had run five times, when four of the five had
  failed at submission. \texttt{sacct} showed four \texttt{FAILED} lines and one
  \texttt{COMPLETED}.
\item it claimed a change had cut the walltime, but the two runs had been given
  different inputs. I re-ran both with the fixed input of
  Section~\ref{sec:lci}.
\item it used an \texttt{sacct} field that does not exist. The command errored.
  The field is in no page of its documentation.
\end{enumerate}

\fi

\end{document}
